% ===================================================================
%  نقشہ نمبر 2 — انسانی جسم۔ — چھٹی۔ — کھیڈاں۔
%  Traduction 2026 d'apres texto/io, controlee sur texto/fr, texto/en
%  et texto/hi. Consigne : voir texto/pnb/10-naqsha-01.tex.
% ===================================================================

%%K t02-apar-1 apar
\VUsaut{4.0mm}
\VUtitre{15.0pt}{60}{نقشہ نمبر 2}
\VUsaut{2.0mm}
\VUtitre{11.4pt}{0}{انسانی جسم۔ --- چھٹی۔}
\VUtitre{11.4pt}{0}{کھیڈاں۔}

%%K t02-tit-0 sub
\VUtitre{13.2pt}{0}{انسانی جسم۔}
\VUsaut{2.4mm}
\VUcentre{10.2pt}{0}{\textit{(قدرتی علماں دا اک سوکھا سبق۔)}}
\VUsaut{3.0mm}

%%K t02-01-1 p
1. --- انسانی جسم دے مکھ حصے ایہہ نیں: \VUgras{سر}، \VUgras{دھڑ} تے
\VUgras{عضو}۔

%%K t02-tit-1 sub
\VUpk{10.2pt}{40}{سر۔}
\VUsaut{3.0mm}

%%K t02-02-1 p
2. --- سر دے دو حصے نیں: \VUgras{کھوپڑی}، جیدے وچ \VUgras{دماغ} رہندا اے تے
جیہنوں \VUgras{وال} ڈھکے رکھدے نیں، تے \VUgras{مونہہ}۔

%%K t02-03-1 p
3. --- ساڈی تصویر وچ نہ کھوپڑی وکھائی دیندی اے نہ دماغ؛ پر مونہہ دے مکھ حصے
سانوں بہت صاف وکھائی دیندے نیں۔

%%K t02-04-1 p
4. --- سب توں پہلاں ایہہ \VUgras{متھا} اے، جیہڑا زوردار عقل تے ہمیشہ کم کردے
دل دا پتہ دیندا اے۔ \VUgras{پٹھ} بہت کھلے نیں۔ \VUgras{اکھ} چنچل تے چوڑی کھلی
اے۔ اک سنی \VUgras{ابرو} تے لمیاں \VUgras{پلکاں} اوہدی رکھیا کردیاں نیں۔

%%K t02-05-1 p
5. --- ایتھے \VUgras{نک} تے \VUgras{نتھنے} وی نیں۔ نک نال اسیں سنگھدے تے ساہ
لیندے آں۔ نک دے دوہاں پاسیں \VUgras{گلھ} نیں، تے اوہناں توں پرے \VUgras{کن}۔
مونہہ دا ہیٹھلا حصہ \VUgras{مونہہ دی گپھا} تے \VUgras{ٹھوڈی} نال بݨدا اے۔

%%K t02-06-1 p
6. --- اوہ سانوں چنگی طرحاں وکھائی نئیں دیندے، کیوں جے \VUgras{مُچھاں} تے
\VUgras{داہڑی} اوہناں نوں لکائی رکھدیاں نیں۔ پر اسیں جاݨنے آں پئی مونہہ وچ دو
\VUgras{بُلھ}، \VUgras{اُتلا جبڑا} تے \VUgras{ہیٹھلا جبڑا} اپݨے
\VUgras{دنداں} سمیت، تے \VUgras{جیبھ} ہوندی اے۔ مونہہ نال اسیں بولنے تے کھانے
آں۔ جیبھ نال اسیں اپݨے کھاݨے دا سواد لیننے آں۔

%%K t02-07-1 p
7. --- اکھ، کن، نک تے مونہہ، انگلاں سمیت، پنجاں حساں دے عضو نیں:
\VUgras{ویکھݨ}، \VUgras{سُݨن}، \VUgras{سنگھݨ}، \VUgras{سواد} تے \VUgras{چھوہ}۔

%%K t02-08-1 p
8. --- ایس طرحاں سر انسانی جسم دے سب توں دلچسپ حصیاں وچوں اک اے۔

%%K t02-tit-2 sub
\VUsaut{4.4mm}
\VUpk{10.2pt}{40}{دھڑ۔}
\VUsaut{3.0mm}

%%K t02-09-1 p
9. --- \VUgras{گردن} سر نوں دھڑ نال جوڑدی اے۔ اوہدے وچ \VUgras{گلا} تے
\VUgras{ڈوڈھی} آندے نیں۔

%%K t02-10-1 p
10. --- دھڑ وچ وی بہت سارے لوڑیندے عضو رہندے نیں۔ \VUgras{چھاتی} وچ
\VUgras{پھیپھڑے} نیں، جیہناں نال اسیں ساہ لیننے آں، تے \VUgras{دل}، جیہڑا
جیوݨ دا کیندر اے۔ جدوں دل دھڑکݨوں ہٹ جاندا اے، جیو مر جاندا اے۔

%%K t02-11-1 p
11. --- \VUgras{پسلیاں} چھاتی نوں سنبھالدیاں نیں۔

%%K t02-12-1 p
12. --- دھڑ دے دوجے حصے نیں \VUgras{پاسے}، \VUgras{پِٹھ}، \VUgras{مونڈھے} تے
\VUgras{ٹِڈھ}۔ اسیں سوݨ لئی پاسے یا پِٹھ دے بھار لیٹنے آں، تے بھار مونڈھے اُتے
چُکنے آں۔

%%K t02-13-1 p
13. --- ٹِڈھ وچ \VUgras{معدہ} تے \VUgras{آندراں} نیں۔ اوہ سانوں کھاݨا ہضم کرن
وچ کم آندیاں نیں۔

%%K t02-tit-3 sub
\VUsaut{4.4mm}
\VUpk{10.2pt}{40}{عضو۔}
\VUsaut{3.0mm}

%%K t02-14-1 p
14. --- ساڈے چار \VUgras{عضو} نیں: دو \VUgras{بانہواں} تے دو \VUgras{لتاں}۔
بانہواں نال اسیں چیزاں لیننے، پھڑنے، چُکنے تے کٹھیاں کرنے آں؛ لتاں نال اسیں
کھڑنے، ٹُرنے تے دوڑنے آں۔

%%K t02-15-1 p
15. --- \VUgras{مونڈھا} بانہہ نوں جسم نال جوڑدا اے۔ اک بہت مضبوط
\VUgras{پٹھا}، دو سِرا، بانہہ نوں ہلاندا اے، جیہنوں \VUgras{ارنڈی}
\VUgras{کلائی والے حصے} نال جوڑدی اے۔ \VUgras{گٹ} \VUgras{ہتھ} دا جوڑ بݨاندا
اے، جیہڑا \VUgras{انگلاں} تے \VUgras{نوہواں} اُتے مُکدا اے۔ ہتھ اُتے سانوں اک
\VUgras{رگ} (تے اک شریان) وکھائی دیندی اے، جیدے وچ \VUgras{لہو} وگدا اے۔

%%K t02-16-1 p
16. --- لت دے حصے ایہہ نیں: \VUgras{پٹ}، \VUgras{گوڈا} تے
\VUgras{گوڈے دا پچھلا پاسا}، \VUgras{پِنّی}، جیدا \VUgras{مانس}
\VUgras{چمڑی} نال ڈھکیا اے، \VUgras{گٹا} تے \VUgras{پیر}۔ لت دیاں
\VUgras{ہڈیاں} بہت موٹیاں تے بہت پکیاں نیں۔ پیر دے سرے اُتے
\VUgras{پیر دیاں انگلاں} نیں۔ دوجے حصے \VUgras{تلی} تے \VUgras{اڈی} نیں۔

%%K t02-17-1 p
17. --- انسانی جسم دی تقریباً پوری وضاحت ایہو ای اے۔

%%K t02-17-2 p
\textit{ٹُوک۔} --- \textit{« میرا ... دُکھدا اے (میرا سر، وغیرہ) » جملے دی مدد
نال ماسٹر ایہہ ساری لفظاولی چنگی یا مندی} \VUgras{جسمانی صحت} \textit{دے پکھوں
مُڑ دہرا سکدے نیں۔}

%%K t02-tit-4 sub
\VUsaut{5.0mm}
\VUpk{10.2pt}{40}{ورزش۔}
\VUsaut{2.4mm}
\VUcentre{10.2pt}{0}{\textit{(اک طالب علم دی زبانی۔)}}
\VUsaut{3.0mm}

%%K t02-18-1 p
18. --- لچکدار، پھرتیلے تے تندرست رہݨ لئی سانوں اپݨیاں لتاں تے بانہواں توں کم
لیݨا چاہیدا اے۔ ایسے لئی اسیں کھیڈاں کھیڈنے آں تے \VUgras{ورزش} کرنے آں۔

%%K t02-19-1 p
19. --- سارا ہفتہ ایہہ دوویں مشقاں سکول دے وِہڑے وچ ہوندیاں نیں (ویکھو نقشہ
نمبر 1)۔ پر ویروار تے اتوار نوں اسیں اک وڈے میدان وچ جانے آں، جِتھے دوڑن تے
موج کرن دی تھاں لبھدی اے۔

%%K t02-20-1 p
20. --- ساڈے ماسٹر تے ساڈے ماں پیو کدے کدے نال آندے نیں؛ پر مشق
\VUgras{ورزش دے ماسٹر} \textsuperscript{(12)} ای کراندے نیں۔

%%K t02-21-1 p
21. --- میدان وچ، وڑن دے کھبے پاسے، \VUgras{ورزش دے سامان}
\textsuperscript{(1)} کھڑے نیں: اسیں \VUgras{پوڑی} \textsuperscript{(2)} اُتے
چڑھنے آں، \VUgras{چھلیاں} \textsuperscript{(3)} تے \VUgras{جھوٹے دے ڈنڈے}
\textsuperscript{(4)} اُتے پُٹھے جھوٹے لیننے آں، \VUgras{رسے دی پوڑی}
\textsuperscript{(5)} یا \VUgras{گنڈھاں والے رسے} \textsuperscript{(6)} دے سرے تیکر
ہتھو ہتھی چڑھ جانے آں۔

%%K t02-22-1 p
22. --- ایس دوران ساڈے ساتھی \VUgras{لکڑ دے گھوڑے} \textsuperscript{(7)} یا
\VUgras{ہم رُخ ڈنڈیاں} \textsuperscript{(11)} دے اُتوں ٹپدے نیں۔ کجھ
\VUgras{مکے بازی} \textsuperscript{(8)} کردے نیں یا نقاب تے \VUgras{فوائل} لے کے
\VUgras{تلوار بازی} \textsuperscript{(9)} کردے نیں۔ کجھ چھیکر \VUgras{ڈمبل}
\textsuperscript{(10)} چُکدے نیں۔

%%K t02-tit-5 sub
\VUsaut{4.6mm}
\VUpk{10.2pt}{40}{کھیڈاں۔}
\VUsaut{3.0mm}

%%K t02-23-1 p
23. --- ورزش کرن والیاں دے چوہاں پاسیں مُنڈے تے کُڑیاں کئی طرحاں دیاں
\VUgras{کھیڈاں} کھیڈ رہے نیں۔ کُڑیاں اپݨے \VUgras{گھیرے} \textsuperscript{(14)}
نال دل پرچاندیاں نیں؛ اوہ \VUgras{رسی ٹپدیاں} \textsuperscript{(13)} نیں، یا
اپݨے بھراواں تے چچیرے بھراواں نوں \VUgras{کروکے} \textsuperscript{(15)} دی کھیڈ
وچ سدّدیاں نیں۔ جدوں مخالف سوچدا اے پئی اوہ سوکھا ای گھیرے وچوں نکل جاوے گا،
تدے اوہ اپݨے \VUgras{لکڑ دے مُنگلی} نال اوہدی \VUgras{گیند} پرے مار دیندیاں
نیں، تے ایس نال اوہناں نوں بہت مزہ آندا اے!

%%K t02-24-1 p
24. --- مُنڈے ودھ زور والیاں کھیڈاں پسند کردے نیں۔ اوہ خوشی نال
\VUgras{نو گولیاں} \textsuperscript{(16)} کھیڈدے نیں، جیہناں نوں اوہ
\VUgras{گیند} \textsuperscript{(17)} نال صاف ڈیگ دیندے نیں۔ اوہ \VUgras{لاٹو}
\textsuperscript{(18)} تے \VUgras{بلبوکے} \textsuperscript{(19)} نوں، ایتھوں تیکر
پئی \VUgras{ڈڈو دی کھیڈ} \textsuperscript{(20)} نوں وی، نکا نئیں سمجھدے۔ پر
اوہناں دیاں پیاریاں کھیڈاں \VUgras{بنٹے} \textsuperscript{(21)} تے
\VUgras{کندھ گیند} \textsuperscript{(22)} نیں، کیوں جے \VUgras{شیشے دے گھر}
\textsuperscript{(26)} دی کندھ ہولی گیند اُچھالݨ لئی ٹھیک بہندی اے۔

%%K t02-25-1 p
25. --- نِکا جان اپݨے \VUgras{بانساں} \textsuperscript{(24)} اُتے چڑھیا بہت
ہوشیار اے تے اپݨا توازن خوب سنبھالدا اے۔ اوہدے کول کجھ ساتھی اوسے شیشے دے گھر
وچ تے \VUgras{واڑ} دے پچھے \VUgras{لُکݨ میٹی} \textsuperscript{(25)} کھیڈ رہے
نیں۔ کجھاں نوں \VUgras{تلی دی کھیڈ} \textsuperscript{(28)} ودھ چنگی لگدی اے، یا
\VUgras{چڑی} \textsuperscript{(29)}، جیہڑی اک \VUgras{بلے} توں دوجے تیکر اُڈدی
رہندی اے۔

%%K t02-noto-1 noto
\VUnotes{143.2mm}{%
(*) نِکیاں دیاں کھیڈاں دے ناں اکثر بالکل قومی ہوندے نیں۔ اسیں اوہناں نوں ایدو
وچ جیویں بݨ سکیا اوویں ناں دِتے نیں — کدے اوس کم توں جیہڑا اوہناں نوں پچھاݨ
دیندا اے، کدے کِسے نہ کِسے دیس دا ناں اپݨا کے، جدوں اوہ بہت رولا پاؤݨ والا نہ
ہووے۔ مثال لئی: \VUgras{پیپے دی کھیڈ} (جرمن تے فرانسیسی) اوہو ای
\VUgras{ڈڈو دی کھیڈ} \textsuperscript{(20)} (انگریزی) اے، جیدے وچ کھڈاری اپݨیاں
گول ٹکیاں موریاں والے صندوق (پیپے) اُتے مونہہ کھولی کھڑے دھات دے ڈڈو دے مونہہ
وچ سُٹن دا جتن کردے نیں۔ \VUgras{مونڈھا ٹپ} \textsuperscript{(30)}
\VUgras{بکری ٹپ} توں نِرا ایہنے کُو وکھری اے: سر دے اُتوں ٹپݨ لئی کھڈاری اپݨے
ساتھی دے مونڈھیاں اُتے ہتھ رکھدے نیں، جدوں کہ « \VUgras{بکری ٹپ} » وچ اوہ نِری
اوہدی پِٹھ اُتے ہتھ رکھدے نیں۔ --- یا فیر، کجھ کھڈاری جھُک جاندے نیں تے دوجے
اوہناں دی پِٹھ اُتے ہتھ رکھ کے اوہناں دے اُتوں ٹپدے نیں؛ پہلیاں نوں بغیر ڈولے
دھکا جرݨا پیندا اے، دوجیاں نوں بغیر توازن گنوائے ٹپݨا (نقشہ ای ویکھو)۔%
}

%%K t02-26-1 p
26. --- میدان دے وچکار دو رُکھ نیں۔ اوہناں وچوں اک دے تھلے ست اٹھ مُنڈیاں نے
\VUgras{مونڈھا ٹپ} \textsuperscript{(30)} (*) دی کھیڈ جمائی ہوئی اے۔ ایہہ کھیڈ ہر
دیس وچ نئیں جاݨی جاندی؛ پر \VUgras{بکری ٹپ} کیہ اے، ایہہ سارے جاݨدے نیں،
جیہنوں نِکے ایس ویلے نئیں کھیڈ رہے۔

%%K t02-tit-6 sub
\VUsaut{5.0mm}
\VUpk{10.2pt}{40}{کھلی ہوا دیاں کھیڈاں۔}
\VUsaut{3.4mm}

%%K t02-27-1 p
27. --- \VUgras{اکھ مچولی} \textsuperscript{(31)} وچ وی بہت مزہ آندا اے؛ کُڑیاں
تے مُنڈے واری واری اکھاں اُتے \VUgras{پٹی} بنھدے نیں تے اوہناں پھوہڑاں نوں
پھڑدے نیں جیہڑے پھڑے جاݨ دیندے نیں۔

%%K t02-28-1 p
28. --- میدان دے وچکار، ایہہ رہی اک بہت فرانسیسی کھیڈ، بھانویں کجھ کھردری:
ایہہ \VUgras{ریچھ} \textsuperscript{(32)} اے، جیہڑی \VUgras{پتنگ}
\textsuperscript{(33)} دی ٹھری کھیڈ توں، ایتھوں تیکر پئی انگریزی کھیڈ
\VUgras{ٹینس} \textsuperscript{(34)} توں وی، کِتے ودھ رولے والی اے۔

%%K t02-29-1 p
29. --- اوہ چھیکرلی کھیڈ وڈے طالب علماں دا آنند اے۔ ساڈے ماسٹر آپ وی اوہنوں
خوشی نال اپݨاندے نیں۔ اوہدے وچ بڑی ہوشیاری تے پھرتی چاہیدی اے، تے اوہ مینوں
بہت چنگی لگدی اے۔ \VUgras{جھوٹا} \textsuperscript{(35)} میں کُڑیاں لئی چھڈ دینا
واں۔

%%K t02-30-1 p
30. --- اک ہور انگریزی کھیڈ وی مینوں چنگی نئیں لگدی، جیہڑی سوکھی ای خطرناک ہو
سکدی اے۔

%%K t02-30-1b p
میرا مطلب \VUgras{فٹ بال} \textsuperscript{(36)} توں اے، جیہڑی میرے وڈے بھرا دا
آنند اے۔ مینوں ساڈی پرانی کھیڈ \VUgras{قیدی اڈہ} \textsuperscript{(38)} ودھ چنگی
لگدی اے، اوہنی ای صحت والی تے کِتے گھٹ سخت۔

%%K t02-tit-7 sub
\VUsaut{4.6mm}
\VUpk{10.2pt}{40}{سائیکل تے گیند دیاں کھیڈاں۔}
\VUsaut{3.0mm}

%%K t02-31-1 p
31. --- مینوں میدان دے چوہاں پاسیں \VUgras{سائیکل} \textsuperscript{(37)} چلاؤݨا
وی چنگا لگدا اے۔

%%K t02-32-1 p
32. --- اگلے ورھے میں \VUgras{گھوڑ سواری} \textsuperscript{(39)} سکھاں گا۔ گھوڑا
کِنّا سوہݨا لگدا اے جدوں اوہ ٹھمکدا یا سوہݨی دلکی ٹُردا اے، تے سرپٹ رکاوٹاں
ٹپ جاندا اے!

%%K t02-33-1 p
33. --- گرمیاں وچ اسیں \VUgras{تیرݨا} \textsuperscript{(40)} سکھنے آں۔ اسیں دریا
دے صاف ٹھنڈے پاݨی وچ آنند نال چُبھی لانے تے نہانے آں۔ کدے کدے چھیکر اسیں اک
کاغذ دا \VUgras{غبارہ} \textsuperscript{(41)} اُڈانے آں، جیہڑا گرم ہوا نال بھردے
ای سنجیدگی نال اُتے چڑھ جاندا اے۔

%%K t02-34-1 p
34. --- ہور وی بہت ساریاں کھیڈاں نیں، پر اوہناں نوں گِݨدے گِݨدے میں کدے نہ
مُکاں گا، تے ایس نِکی جیہی کہاݨی توں ای وکھائی دیندا اے پئی ساڈے سوہݨے میدان
اُتے ویروار نیرس ہوݨوں بہت دور نیں۔

%%K t02-34-2 p
نوٹ۔ --- \textit{ماسٹر ایس باب نوں سوکھا ای پورا کر سکدے نیں، ہر اہم کھیڈ دی
ودھ تفصیل نال وضاحت کر کے۔}
\VUsaut{6.0mm}
\VUfilet{22mm}
